% Transcription — fonds Alexandre Grothendieck, shelfmark 43, batch 9 (pages 161-180).
% Established from the facsimile archives/batches/43/batch-09.pdf (PDF page k =
% archivists' page 160+k, no cover sheet), cut from a copy of 43.pdf fetched
% through the project's own relay
% (https://grothendieck-relay.onrender.com/source/43.pdf); tiles in
% archives/tiles/43/p161-p180. Per the skill transcribe-grothendieck. The folder
% is 207 pages in eleven batches, transcribed in parallel; no neighbouring
% batch file existed when this one was written, so the notation is fixed from
% these pages alone (and aligned with folders 59 and 9 for G_m, G_a, alpha_p,
% mu_p, W, Hom/Ext underlined).
%
% Pass: Opus 5.5 (claude-opus-5-5), 2026-09-26 — first pass, unchecked against the pages by a human.
%
% Dating: the folder has its own inventory date, carried verbatim; the group
% « Jacobiennes » (43-44) is dated [à partir de 1961-à partir de 1967] and is
% given after it, as the skill's group rule asks. No page of this batch
% carries a date.
%
% Copies. The folder title speaks of « copies de notes manuscrites ». In this
% batch, pages 165, 174 and 178 are very faint copies on cream paper (not
% mirror-image show-through: the text reads the right way round) of,
% respectively, pages 166, 175 and 179; as far as the faint strokes allow,
% their layout and wording coincide with those leaves. Pages 161-169 are on
% a grey ground with uniformly grey strokes and look like photocopies of
% pencil or ink notes (page 161 is written across a blue-grey sheet turned
% sideways); pages 171-180 are originals in blue-black ink on « J.A. VELIN
% EXTRA » paper (watermark), page 170 is an orange cover sheet.
%
% Pages skipped: none. Page 170 is a cover (orange card) inscribed in his
% hand « Dualité de Cartier des groupes affines et groupes formels »: it
% takes no \page marker and its words become the section heading, with a
% \note.
%
% Pages summarised: 165, 174, 178 — the faint copies above, each given one
% \page marker and a \note pointing to the leaf it copies, not re-transcribed.
% Page 177 (overlapping ovals classifying the basic groups) is transcribed
% by its labels and prose, the layout of the ovals being described in \note{}s.
%
% His pagination: none visible on these twenty leaves. There is no numbered
% run; page 168 refers to « l'exposé du livre de Serre ... p. 183 », and
% page 169 to « Exposé 7 » of a Chevalley seminar (in the margin).
%
% What the batch is about. Hom and Ext^1 between the basic commutative
% affine groups over a base S of characteristic p: Z/lZ, Z/pZ, alpha_p, mu_p,
% G_m, G_a, Z, the Witt covector group (bold V), an abelian scheme A.
% P. 161 (sheet turned): a first Hom table, cancelled; then « Tableau des
% Hom » with rows and columns Z/lZ, Z/pZ, alpha_p, mu_p, G_m, G_a, Z, V, A
% and a 3x3 matrix of Hom between mu_p, alpha_p, Z/pZ. P. 162: the rules
% that produce the rows (Hom(Z/nZ,G) = _nG, Hom(alpha_p,G), Hom(mu_p,G) as
% kernels on V(omega_G), Hom(V,G) = V(omega_G), « non ramifié » functors)
% and the « cas difficiles » Hom(G_a,G), Hom(G,V). P. 163: a table of Ext^1
% (rows _pG_m = mu_p, _FG_a = W_11, Z/lZ, Z/pZ, G_m, G_a, Ab, Z), with
% Ext^1(G_m,_pA), lim Ext^1(G_m,_nA), Hom(T(A),G), H^1(Ab,O), and « Consulter
% travaux Serre ». P. 164: a coarser table by types (Gamma finite ordinary, U
% infinitesimal, G_m, G_a, Ab simple, Z), « Tableau des Hom », and
% Hom(G_a,_pG_m) via unipotent elements of k[[t]]. Pp. 166-168: Lemma —
% for G, Gamma commutative with Hom_S(G,Gamma) = 0, Ext^1_{S-gr}(G,Gamma) ->
% H^1(G,Gamma_G) is injective with image the primitive classes; proof
% (making a section multiplicative; building the group law on E;
% associativity, unit, symmetry, « voir ... Serre p. 183 »); application:
% for Gamma étale and G with geometrically connected fibres, Ext^1 is
% invariant under S_red -> S (more generally a universal homeomorphism).
% P. 169: example Ext^1_{S-gr}(G_a, Z/pZ) = Gamma(S, O^{p^{-infty}}),
% representable by perfect pro-schemes. Pp. 171-180, under the cover
% « Dualité de Cartier des groupes affines et groupes formels »: groups
% classified by the behaviour of Frobenius F and Verschiebung V (isom,
% épim, nilpotent; tables pp. 171-173), the « 8 types fondamentaux »,
% the dictionary multiplicatif / unipotent / discret / infinitésimal
% (p. 175), the table G_m, Z/lZ, mu_p / G_a, Z/pZ, alpha_p (p. 176),
% diagrams of classes (p. 177, 179), and p. 180: discrete modules of finite
% type over the Dieudonné ring A = lim W_n[F,V]/(V^n, FV - p), categories
% C_1, C_2 = C_1°, Cartier duality C_1' = C_2', and a glued category of
% triples (X, Y, u).
%
% Notation fixed once for the batch: \mathbb{G}_m, \mathbb{G}_a, \alpha_p,
% \mu_p, \mathbb{Z}/\ell\mathbb{Z} (his l is \ell), bold outline V as
% \mathbb{V} (the Witt covector group, as he writes it double-struck), A for
% an abelian scheme, \underline{\mathrm{Hom}}, \underline{\mathrm{Ext}} where
% he underlines, \underline{V}(\omega_G) for his underlined V, \Gamma for his
% curly capital (a discrete / étale group) on pp. 164-169, _pG, _FG_a for
% his left subscripts, W_{11} as written on p. 163. Tables are set as arrays
% where the cells are formulas, as lists where the cells are words.
%
% Readings to check first: the cancelled first table and the pencil margin
% of p. 161; the circled cells (« gros », « ? ») of p. 161; the letter
% written under Hom(alpha_p, ...) on p. 162 (read as a script L); the
% row/column assignments of p. 163 (cells marked ∅, ×, * are his); the
% c') row of p. 164; « prédit » on p. 166; the diagonal margin of p. 167;
% the marginal reference on p. 169; the ovals of p. 177; « cocyclique »
% (?) and the last lines of p. 180.

\documentclass[11pt,a4paper]{article}
% Le préambule commun à toutes les transcriptions du fonds Grothendieck.
%
% Il tient en une idée : **l'incertitude se compose, elle ne se résout pas.**
% Une transcription qui lisse un mot douteux a détruit la seule chose pour
% laquelle on la faisait — un lecteur doit pouvoir savoir, à chaque endroit,
% ce qui était sur le feuillet et ce qui a été ajouté. Les six macros qui
% suivent sont donc l'appareil critique, et non de la décoration.
%
% Les mêmes commandes sont reconnues par `scripts/render.mjs`, qui produit la
% vue de lecture du site. Une macro ajoutée ici doit l'être aussi là-bas,
% faute de quoi le rendu échouera bruyamment — ce qui est voulu : un rendu
% silencieusement faux est pire qu'un rendu absent.

\ProvidesPackage{grothendieck}[2026/08/08 Transcriptions du fonds Grothendieck]

\usepackage{amsmath,amssymb,amsthm}
\usepackage{mathrsfs}
\usepackage{stmaryrd}
% Les diagrammes commutatifs sont partout dans ce fonds : c'est la langue
% même de la géométrie algébrique de Grothendieck, et une transcription qui
% les rendrait en prose ne transcrirait rien.
\usepackage{tikz-cd}
% Français par défaut — c'est la langue des feuillets — mais l'anglais est
% chargé aussi : la traduction et le résumé sont des documents anglais, et la
% typographie française (espaces avant les deux-points) y serait une faute.
% Ces documents basculent par \selectlanguage{english} en tête de corps.
\usepackage[english,french]{babel}
\usepackage{fontspec}
\usepackage{geometry}
\geometry{a4paper,margin=25mm,marginparwidth=28mm}
\usepackage{marginnote}
\usepackage{xcolor}
% ulem, not soul. `soul`'s \st and \ul are built on a character-by-character
% scan that XeLaTeX's font machinery breaks: the document fails with an
% undefined control sequence rather than degrading. ulem does the same two jobs
% and is Unicode-engine safe. `normalem` stops it hijacking \emph, which the
% transcriptions use for Grothendieck's own emphasis.
\usepackage[normalem]{ulem}
\usepackage{hyperref}
\hypersetup{colorlinks=true,linkcolor=black,urlcolor=black,pdfborder={0 0 0}}

% --- Métadonnées du lot -----------------------------------------------------
% Reprises telles quelles par le script de rendu, qui les affiche en tête de
% la vue de lecture. Elles fixent ce que le fichier prétend couvrir : sans
% elles, une transcription flotte sans point d'attache dans un fonds de seize
% mille feuillets.

\newcommand{\folder}[1]{\gdef\grothFolder{#1}}
\newcommand{\batch}[1]{\gdef\grothBatch{#1}}
\newcommand{\leaves}[2]{\gdef\grothFirst{#1}\gdef\grothLast{#2}}
\newcommand{\foldertitle}[1]{\gdef\grothFoldertitle{#1}}
\gdef\grothFolder{?}\gdef\grothBatch{?}\gdef\grothFirst{?}\gdef\grothLast{?}\gdef\grothFoldertitle{}

% --- Le résumé --------------------------------------------------------------
%
% Il ouvre la lecture modernisée, sous le titre « Résumé », et il est composé
% en petit. Ce n'est pas un ornement : c'est le seul passage du document qui
% ne soit pas des mathématiques, et un lecteur doit pouvoir le reconnaître
% avant de l'avoir lu — puis le sauter s'il connaît déjà le sujet, ou s'y
% arrêter s'il ne le connaît pas. Le filet qui le suit marque la reprise du
% corps.

\newenvironment{resume}
  {\small\setlength{\parskip}{0.45em}}
  {\par\medskip\hrule\medskip}

% --- Mots-clés ---------------------------------------------------------------
%
% En anglais, en clôture du résumé de la lecture modernisée : le vocabulaire
% moderne sous lequel chercher ce que les feuillets construisent. Le manifeste
% du site (`npm run manifest`) les extrait pour étiqueter la cote entière —
% cette ligne est la source unique des tags, il n'y a pas de fichier à côté.
\newcommand{\keywords}[1]{\par\smallskip\noindent
  {\footnotesize\textsc{Keywords} --- \textit{#1}}\par}

% --- L'appareil critique ----------------------------------------------------

% Le numéro de feuillet, en marge. C'est l'ancre qui permet de régler un
% désaccord sur une formule en regardant *un* feuillet plutôt qu'une chemise
% de six cents. Le site s'en sert aussi pour tourner les pages du fac-similé
% quand on fait défiler la transcription.
\newcommand{\leaf}[1]{%
  \marginnote{\footnotesize\color{gray!70}\textsf{#1}}%
  \label{leaf:#1}\ignorespaces}

% Illisible : on ne devine pas. Un mot inventé qui se lit comme les autres est
% le pire résultat possible.
\newcommand{\ill}{\textcolor{red!60!black}{[\,\ldots\,]}}

% Lecture douteuse : le mot est donné, et signalé comme incertain.
\newcommand{\uncertain}[1]{\uline{#1}}

% Ajout de l'éditeur — une lettre manquante, un mot sous-entendu.
\newcommand{\add}[1]{\textcolor{blue!50!black}{[#1]}}

% Biffé par Grothendieck lui-même. Ce qu'il a rayé fait partie du document :
% une rature dit souvent où la pensée a changé de direction.
\newcommand{\struck}[1]{\sout{#1}}

% Note du transcripteur, sur l'état matériel du feuillet.
\newcommand{\note}[1]{\par\smallskip{\small\color{gray!60!black}\textit{[#1]}}\par\smallskip}

% Note marginale de Grothendieck, distincte du corps du texte.
\newcommand{\marginal}[1]{\par\smallskip\noindent\hspace{1em}%
  {\small\color{gray!60!black}#1}\par\smallskip}

% --- Titre ------------------------------------------------------------------

\AtBeginDocument{%
  \begin{center}
    {\large\bfseries Fonds Alexandre Grothendieck --- cote n\textsuperscript{o}~\grothFolder}\\[2pt]
    {\itshape\grothFoldertitle}\\[4pt]
    {\small feuillets \grothFirst--\grothLast\ (lot \grothBatch)}\\[2pt]
    {\footnotesize\color{gray!60!black}Transcription établie d'après le fac-similé numérisé
     par l'Université de Montpellier.\\
     Les lectures incertaines sont soulignées, les illisibles notées [\,\ldots\,],
     les ajouts entre crochets.}
  \end{center}
  \vspace{1em}}

\endinput


\folder{43}
\batch{9}
\pages{161}{180}
\dating{[à partir de 1961-vers 1968]}
\watermark{Édition de démonstration}
\foldertitle{Dualité des groupes. Jacobiennes généralisées etc : notes et copies de notes manuscrites (s.d.)}

\begin{document}

\page{161}
\note{le feuillet est écrit en travers (tourné d'un quart de tour) ; il
semble être une photocopie. Dans la marge (au crayon, verticalement), des
mots presque effacés : \ill{} \ill{} $+\,\mathrm{Ext}^1$ \ill{} \ill{}, \ill{}}

\struck{Ind pro}\qquad \emph{Ext}${}^1(F,G)$

\note{un premier tableau, biffé par de longues diagonales :}
\[
\begin{array}{c|c|c|c|c|c|c|c|c|c}
\underline{\mathrm{Hom}}(F,G) & \mathbb{Z}/\ell\mathbb{Z} & \mathbb{Z}/p & \mu_p & \alpha_p & \mathbb{G}_m & \mathbb{G}_a & \mathbb{V} & \mathbb{Z} & A \\ \hline
\mathbb{Z}/\ell\mathbb{Z} & \mathbb{Z}/\ell\mathbb{Z} & & 0 & 0 & \mu_\ell & 0 & 0 & 0 & {}_\ell A \\
\mu_p & 0 & & \mathbb{Z}/p\mathbb{Z} & & & & & &
\end{array}
\]
\note{le $\mathbb{Z}$ de l'avant-dernière colonne est lui-même biffé ; la
colonne $\mathbb{Z}/p$ est ajoutée, serrée, après $\mathbb{Z}/\ell\mathbb{Z}$}

\emph{Tableau des Hom}

\struck{$\mu_p$ $\alpha_p$ $\mathbb{Z}/p\mathbb{Z}$ $\mathbb{Z}/\ell\mathbb{Z}$ Tableau des Ext}

\[
\begin{array}{c|c|c|c|c|c|c|c|c|c}
 & \mathbb{Z}/\ell\mathbb{Z} & \mathbb{Z}/p\mathbb{Z} & \alpha_p & \mu_p & \mathbb{G}_m & \mathbb{G}_a & \mathbb{Z} & \mathbb{V} & A \\ \hline
\mathbb{Z}/\ell\mathbb{Z} & \mathbb{Z}/\ell\mathbb{Z} & 0 & 0 & 0 & \mu_\ell & 0 & 0 & 0 & {}_\ell A \\
\mathbb{Z}/p\mathbb{Z} & 0 & \mathbb{Z}/p\mathbb{Z} & \alpha_p & \mu_p & \mu_p & \mathbb{G}_a & 0 & \mathbb{V} & {}_p A \\
\alpha_p & 0 & 0 & \mathbb{G}_a & \alpha_p & \alpha_p & \mathbb{G}_a & 0 & (\mathrm{gros}) & ? \\
\mu_p & 0 & 0 & 0 & \mathbb{Z}/p\mathbb{Z} & \mathbb{Z}/p\mathbb{Z} & 0 & 0 & 0 & ? \\
\mathbb{G}_m & 0 & 0 & 0 & 0 & \mathbb{Z} & 0 & 0 & 0 & 0 \\
\mathbb{G}_a & 0 & 0 & (\mathrm{gros}) & \mathbb{V} & \mathbb{V} & ? & 0 & (\mathbb{G}_a ?) & \mathrm{gros} \\
\mathbb{Z} & \mathbb{Z}/\ell\mathbb{Z} & \mathbb{Z}/p\mathbb{Z} & \alpha_p & \mu_p & \mathbb{G}_m & \mathbb{G}_a & \mathbb{Z} & \mathbb{V} & A \\
\mathbb{V} & 0 & 0 & \mathbb{G}_a & \mathbb{G}_a & \mathbb{G}_a & \mathbb{G}_a ? & 0 & (\mathrm{Imm}) & \underline{V}(\omega_A) \\
A & 0 & 0 & 0 & 0 & 0 & 0 & 0 & 0 ? & \mathbb{Z}^n
\end{array}
\]
\note{les lignes sont le premier argument, les colonnes le second : la
ligne $\mathbb{Z}$ redonne l'en-tête. « gros », « Imm », « $\mathbb{G}_a$ ? »
sont entourés. Dans la colonne $A$, les cases des lignes $\alpha_p$ et $\mu_p$
sont couvertes de hachures, avec les restes lisibles « ${}_{1-F}A$ ? » et
« $({}_{F}A)$ ? » ; la case $(\mathbb{G}_a,\mathbb{G}_a)$ porte un mot biffé
\ill{} et, en dessous, « Imm ? » ; la case $(\mu_p,\mathbb{V})$ porte un
« 0 » surchargé. Une accolade, à gauche des lignes $\alpha_p$ à $\mathbb{G}_a$,
porte « \uncertain{conn} »}

\[
\begin{pmatrix}
\mathbb{Z}/p\mathbb{Z} & 0 & 0 \\
\alpha_p & \mathbb{G}_a & 0 \\
\mu_p & \alpha_p & \mathbb{Z}/p\mathbb{Z}
\end{pmatrix}
\]

\emph{Hom}$(\alpha_p, \mathbb{G}_m)$ \qquad \ill{} ??? \ill{}
\note{la dernière inscription, à droite, est entourée d'un ovale au crayon}

\page{162}
\[ \underline{\mathrm{Hom}}(G,H) = 0 \]
si $G$ à fibres géom. connexes, $H$ non ramifié sur $S$.
\[ \underline{\mathrm{Hom}}(\mathbb{Z}/n\mathbb{Z}, G) \simeq {}_nG \]
d'où lignes 1, 2
\[ \underline{\mathrm{Hom}}(\mathbb{Z}, G) \simeq G \]
d'où ligne 7
\[
\underline{\mathrm{Hom}}(\alpha_p, G) \simeq \mathrm{Ker}\bigl(\underline{V}(\omega_G) \xrightarrow{\ \text{puiss. $p$-ième}\ } \underline{V}(\omega_G)\bigr)
\;\simeq\; \underline{\mathrm{Hom}}(\alpha_p, \underline{\mathfrak{L}}^{(1)}_G)
\]
\note{sous $\underline{V}(\omega_G)$ est écrit « $\underline{\mathfrak{L}}_G$
schéma de Lie de $G$ » ; la lettre gothique est une lecture incertaine}
d'où ligne 3 \uncertain{moyennant} l'\ill{} \ill{} ----
\[
\underline{\mathrm{Hom}}(\mu_p, G) \simeq \mathrm{Ker}\bigl(\underline{V}(\omega_G) \xrightarrow{\ \text{id $-$ puiss. $p$-ième}\ } \underline{V}(\omega_G)\bigr)
\]
d'où ligne 4

\emph{Hom}$(\mathbb{G}_m, G)$ est un foncteur « non ramifié », d'où ligne 5

\emph{Hom}$(A, G)$ est un foncteur « non ramifié », d'où (partiellement) ligne 6.

\emph{Hom}$(\mathbb{V}, G) \simeq \underline{V}(\omega_G)$ \quad d'où ligne 8
\note{les numéros de ligne renvoient au tableau des Hom de la page 161, lu
dans l'ordre $\mathbb{Z}/\ell\mathbb{Z}$, $\mathbb{Z}/p\mathbb{Z}$,
$\alpha_p$, $\mu_p$, $\mathbb{G}_m$, $\mathbb{G}_a$, $\mathbb{Z}$,
$\mathbb{V}$ ; la ligne 6 attribuée à $A$ ne coïncide pas avec cet ordre}

Cas difficiles

a) \emph{Hom}$(\mathbb{G}_a, G)$, \quad $G = \alpha_p, \mathbb{G}_a, \mathbb{V}$, $/\!/A$

b) \emph{Hom}$(G, \mathbb{V})$, \quad $G = \alpha_p, \mathbb{G}_a, \mathbb{V}$, $/\!/A$

(\uncertain{indiqué} \ill{} par $G = \mathbb{G}_a, \mathbb{V}, \alpha_p$)

Il y a essentiellement \struck{deux} \add{quatre} types « délicats » -----
\note{« quatre » est écrit au-dessus du mot biffé}

\page{163}
\[
\begin{array}{c|c|c|c|c|c|c|c|c}
 & \mu_p & W_{11} & \mathbb{Z}/\ell\mathbb{Z} & \mathbb{Z}/p\mathbb{Z} & \mathbb{G}_m & \mathbb{G}_a & Ab & \mathbb{Z} \\ \hline
{}_p\mathbb{G}_m = \mu_p & \mathbb{Z}/p & 0 & 0\ \ast & \ast\ 0 & ? & 0 & [a] & 0\ \ast \\
{}_F\mathbb{G}_a = W_{11} & 0 & \mathbb{G}_a^2\ (?) & 0\ \ast & \ast\ 0 & 0 & (?) & [b] & 0\ \ast \\
\mathbb{Z}/\ell\mathbb{Z} & 0\ \times & 0\ \times & \mathbb{Z}/\ell\mathbb{Z} & 0 & \emptyset\ 0 & \emptyset\ 0 & \emptyset\ 0 & \emptyset\ \mathbb{Z}/\ell \\
\mathbb{Z}/p\mathbb{Z} & 0\ \ast & 0\ \times & 0 & \mathbb{Z}/p\mathbb{Z} & \emptyset\ 0 & \mathbb{G}_a & \emptyset\ 0 & \emptyset\ \mathbb{Z}/p\mathbb{Z} \\
\mathbb{G}_m & \mathbb{Z}/p\mathbb{Z} & 0 & [c] & 0 & 0 & 0 & [d] & \\
\mathbb{G}_a & 0 & (?) & 0\ 0 & (?) & 0 & (?) & [e] & \\
Ab & [f] & \longleftarrow & \longrightarrow & & (Ab)' & H^1(Ab, \underline{O}) & & \\
\mathbb{Z} & 0\ \ast & 0\ \ast & 0\ \ast & 0\ \ast & 0\ \ast & 0\ \ast & 0\ \ast & 0\ \ast
\end{array}
\]
\note{tableau des $\mathrm{Ext}^1$ (lignes : premier argument). Les
renvois [a] à [f] sont de l'éditeur, faute de place : [a]
$\mathrm{Ext}^1(\mathbb{G}_m, {}_pA)$ ; [b] $\mathrm{Ext}^1(\mathbb{G}_a,
{}_FA)$ ; [c] $\mathrm{Hom}(T(\mathbb{G}_m), \mathbb{Z}/\ell\mathbb{Z})$ ;
[d] $\varinjlim \mathrm{Ext}^1(\mathbb{G}_m, {}_nA)$ (« lim » souligné
deux fois) ; [e] $\varinjlim \mathrm{Ext}^1(\mathbb{G}_a, {}_nA)$ ; [f]
$\mathrm{Hom}(T(A), G)$. La case
$(\mu_p,\mu_p)$ porte en surcharge « \ill{} dualité » ; la case
$(\mu_p,\mathbb{G}_m)$ est surchargée, illisible ; dans la case
$(\mathbb{G}_m,\mathbb{Z}/\ell\mathbb{Z})$, sous le Hom, un
« $\mathbb{Z}/\ell\mathbb{Z}$ » partiellement biffé ; la case $(Ab, Ab)$
porte un mot biffé \ill{} ; une grande accolade occupe la colonne
$\mathbb{Z}$ en face des lignes $\mathbb{G}_m$ à $Ab$ ; dans la ligne
$Ab$, « $\mathrm{Hom}(T(A), G)$ » est écrit sur une flèche qui couvre les
quatre premières colonnes. Les signes $\emptyset$,
$\times$, $\ast$ sont les siens et renvoient aux remarques ci-dessous. Deux
crochets à gauche groupent les lignes ${}_p\mathbb{G}_m$ à
$\mathbb{Z}/p\mathbb{Z}$ et $\mathbb{G}_m$ à $\mathbb{Z}$}

$\emptyset$ \ $\mathrm{Ext}^1(\mathbb{Z}/n\mathbb{Z}, G) = {}^nG$,
\uncertain{$\Gamma$ discret fini}, $G$ \uncertain{dualisable}

$\times$ \ \struck{\ill{}} $\mathrm{Ext}^1(\Gamma, G) = 0$ si \ill{} dans
\struck{\ill{}} [\emph{discret}], l'un \struck{\ill{}}, l'autre
\uncertain{infinitésimal}, \ill{} \ill{} ext. \ill{} \ill{}

\marginal{$\mathrm{Ext}^1$ dans la cat. \emph{géométrique}}
\note{dans un nuage tracé à droite des deux remarques précédentes}

\struck{N.B. $\mathrm{Ext}^2(\mathbb{G}_m, \mathbb{G}_m) \neq 0$ \ill{} \ill{}}
\note{encadré et biffé par des boucles}

(?) \emph{Consulter travaux Serre} \uncertain{pour} \ill{} déterminations explicites
\note{le « ? » est entouré}

(Il \ill{} le dit qu'il y a \uncertain{lieu} de \uncertain{considérer} des
cat. de $G$, $\mathbb{G}_m$ \ill{} $Ab$ \ill{} $\mathbb{Z}$ ; on trouve
\ill{} \ill{}, dont 0

\page{164}
\note{tableau à six colonnes $\Gamma$, $U$, $\mathbb{G}_m$, $\mathbb{G}_a$,
$Ab$ \uncertain{simple}, $\mathbb{Z}$ ; la case d'en-tête de la première
colonne est un trait. Les lignes, repérées à gauche par des lettres, sont
données ici une à une, case par case dans l'ordre des colonnes}

\begin{itemize}
\item b), b') $\Gamma \in \lbrace \mathbb{Z}/\ell\mathbb{Z}, \mathbb{Z}/p\mathbb{Z} \rbrace$ :
groupe fini ordinaire \ill{} $\mathrm{Hom}(\Gamma, \Gamma')$ ;
0 ou \struck{\ill{}} groupe radiciel simple ;
groupe simple (ext. ou radiciel) ;
0 ou $\mathbb{G}_a$ ;
gr. \ill{} \ill{} mixte ;
0.
\item a), a') $U \in \lbrace {}_p\mathbb{G}_m = D(\mathbb{Z}/p\mathbb{Z}),\ {}_F\mathbb{G}_a \rbrace$ :
0 ;
$\mathbb{Z}/p\mathbb{Z}$, \struck{$\mathbb{Z}$} ? pour ${}_p\mathbb{G}_m$, et
${}_F\mathbb{G}_a$, $\mathbb{G}_a$ pour ${}_F\mathbb{G}_a$ ;
$\mathbb{Z}/p\mathbb{Z}$, ${}_F\mathbb{G}_a$ ;
0, $\mathbb{G}_a$ ;
groupe $(\mathbb{Z}/p\mathbb{Z})^n$, ${}_F\mathbb{G}_a^m \times \mathbb{G}_a^n$ ;
0.
\item c) $\mathbb{G}_m$ :
0 ; $\times$ 0 ; $\mathbb{Z}$ ; 0 ; 0 ; 0.
\item c') $\mathbb{G}_a$ :
0 ;
$C W_0$ ${}\times^{(\alpha)}$ (\uncertain{Hom conoyaux} \ill{} \ill{} fini \ill{})
et ${}^{(\beta)}$ \ill{} \uncertain{vecteurs} \ill{} \ill{} ;
radiciel \ill{} $= \mathrm{Hom}(\mathbb{G}_a, {}_p\mathbb{G}_m)$ ;
$\mathbb{V}$ $\times$ ;
radiciel $= \mathrm{Hom}(\mathbb{G}_a, U)$ ;
0.
\item d) $Ab$ \uncertain{simple} :
0 ; $\times$ 0 ; 0 ; 0 ;
groupe discret \ill{} de type fini sans torsion ;
0.
\item e) $\mathbb{Z}$ groupe ordinaire :
$\Gamma$ ; $U$ ; $\mathbb{G}_m$ ; $\mathbb{G}_a$ ; $Ab$ \uncertain{simple} ; $\mathbb{Z}$.
\end{itemize}
\note{dans la case $(U, U)$ un petit tableau $2 \times 2$ séparé par des
pointillés ; les croix $\times$ sont à l'encre rouge ; « 0 ou » est ajouté
au-dessus de la case $(\Gamma, U)$}

\emph{Tableau des Hom}
\note{encadré}

\emph{Hom}(gr. connexe, gr. discret ordinaire) $= 0$

\emph{Hom}($\mathbb{G}_m$, gr. \ill{} \ill{} \ill{} $\mathbb{G}_m$) \struck{$\neq$} $= 0$

\emph{Hom}($Ab$, gr. linéaire) $= 0$

$\mathbb{S}_\infty \longrightarrow \mathbb{G}_a$ \quad hom. d'anneaux de
$\mathbb{S}_\infty$ dans $\mathbb{G}_a$
\note{la lettre ajourée est lue $\mathbb{S}$ ; lecture incertaine}

\struck{(\ill{})(\ill{})}

$\underline{\mathrm{Hom}}(\mathbb{G}_a, {}_p\mathbb{G}_m)$
\[
k[[t]] \longleftarrow k[t]/(1-t)^p
\]
\uncertain{éléments} unipotents en $k[[t]]$, i.e. de la forme
$\sum a_i t^i$, les $a_i^p = 0$
\struck{$\Delta(\sum a_i t^i) =$}
\[
1 + a_1 t + a_2 t^2 + \cdots + a_n t^n \qquad a_i^p = 0
\]
\[
P(t+t') = P(t)\,P(t')
\]

\marginal{groupe (\ill{} algébrique) ; $\alpha$) L'\ill{} \struck{\ill{}} $D(\mathbb{G}_a)$ de Cartier ---- \ill{} un
gr. \emph{radiciel infini} ; $\beta$) \ill{} un groupe $k$-simple infini sur $k$.}
\note{colonne de droite, séparée par un trait ; elle se rapporte aux
renvois $(\alpha)$, $(\beta)$ de la ligne c') du tableau}

\page{165}
\note{copie très pâle du feuillet suivant (page 166) : même lemme, même
disposition (« Lemme. Soient $G$, $\Gamma$ \ldots », la formule (*), les
accolades de la marge gauche) ; non retranscrite, voir page 166}

\page{166}
\marginal{1° Tout \uncertain{hom.} $G \times_S G \to \Gamma$ qui est nul sur
$e \times_S G$ et sur $G \times_S e$ est nul ;
2° Si tout \uncertain{hom.} $(G/e)^n \to \Gamma$ \ill{} \ill{} \ill{} \ill{}
$S \to \Gamma$ \struck{alors (**)} (pour $n = 1, 2, 3$) alors}
\note{deux conditions écrites verticalement dans la marge gauche, reliées
au texte par des accolades}

\emph{Lemme}. Soient $G$, $\Gamma$ deux schémas en groupes
\marginal{commutatifs} sur $S$, tels que
\struck{$\underline{\mathrm{Hom}}_S(G,\Gamma) = 0$} [par exemple $G$ à fibres
géométriquement connexes, et $\Gamma$ non ramifié sur $S$], $\Gamma$ fid.
plat sur $S$ pour simplifier. Alors
\note{la condition $\underline{\mathrm{Hom}}_S(G,\Gamma) = 0$ est barrée
d'un trait ondulé ; au-dessus, « Hom » ; en dessous,
« \uncertain{prédit} » ; « commutatifs » est dans une bulle en marge}
\[
(*)\qquad \mathrm{Ext}^1_{S\text{-gr}}(G, \Gamma) \longrightarrow H^1(G, \Gamma_G)
\]
\note{sous $H^1(G,\Gamma_G)$, une accolade : « classes de fibrés principaux
homogènes »}
est \emph{injectif}, et \add{a} pour image l'ensemble des éléments
\emph{primitifs} de $H^1(G, \Gamma_G)$ [i.e. ceux dont l'image dans
$H^1(S,\Gamma)$ sont nuls, et satisfaisant
$\mu^{*}(\xi) - \mathrm{pr}_1^{*}(\xi) - \mathrm{pr}_2^{*}(\xi) = 0$]
\note{le premier terme de cette relation est de lecture incertaine}

\emph{Dém}. Soit
\[
0 \to \Gamma \xrightarrow{\ \alpha\ } E \xrightarrow{\ \beta\ } G \to 0
\]
une extension de $S$-groupes [$E \to G$ fid. plat], supposons \ill{} \ill{}
$E$ comme fibré principal \emph{trivial} sur $G$, i.e. qu'il existe une
section $\varphi : G \to E$. Soit $e_G : S \to G$ la section unité, d'où une
section $\varphi \circ e_G = \varphi(e_G)$ de $E$ sur $S$, dont l'image par
$\beta$ est la section unité $e_G$, donc qui est de la forme $\alpha(u)$,
$u \in \Gamma(S)$. Corrigeant

\page{167}
$\varphi$ par $u$, on peut supposer que $\varphi$ transforme section unité
en section unité. Prouvons que $\varphi$ est \emph{multiplicative}. En effet,
$\pi' \circ (\varphi \times_S \varphi) - \varphi \circ \pi$

\begin{tikzcd}
G \times G \arrow[r, "\varphi \times \varphi"] \arrow[d, "\pi"'] & E \times E \arrow[d, "\pi'"] \\
G \arrow[r, "\varphi"] & E
\end{tikzcd}

est un $S$-morphisme $G \times G \to \Gamma$, qui est nul sur
$G \times_S e$ et $e \times_S G$.
\note{le $\Gamma$ de $G \times G \to \Gamma$ est écrit en surcharge, très
noirci}

Donc \uncertain{ceci prouve} l'injectivité ----.

Le fait que l'image de
$\mathrm{Ext}^1_{S\text{-gr}}(G,\Gamma) \to H^1(G,\Gamma)$ soit les primitifs
\uncertain{indiqués}, est vrai sans hyp. sur $G$, $\Gamma$, compte tenu de la
construction des fibrés associés. Pour la réciproque,
\add{sous les conditions dites 2°} on note que \struck{\ill{}} $E$ est fibré
pr. homogène sur $G$ \struck{\ill{}} de \add{classe} primitive, il existe
\struck{\ill{}} une section $e_E$ de $E/S$ au dessus de celle de $G/S$.

\begin{tikzcd}[column sep=small]
\Gamma \times \Gamma & E \times E \arrow[r] \arrow[d] & G \times G \arrow[d] \\
\Gamma & E \arrow[r] & G
\end{tikzcd}

De plus, il existe un morphisme $E \times E \to E$ \ill{}
$\Gamma \times \Gamma \to \Gamma$ et rendant commutatif

\begin{tikzcd}
E \times E \arrow[r] \arrow[d] & G \times G \arrow[d] \\
E \arrow[r] & G
\end{tikzcd}

Il y a \ill{} multiplications \uncertain{possibles} dans $E$
\add{$\underline{\mathrm{Hom}}(G\times G, \Gamma) \simeq \underline{\mathrm{Hom}}(G,\Gamma)$ ?}
\ill{} \ill{} \ill{} \ill{} \ill{} \ill{}, \ill{} la multiplication
\struck{\ill{}} (unique) qui transforme $(e_E, e_E)$ en $e_E$.
\note{l'ajout entre crochets, entouré d'une boucle, est de lecture
incertaine}

\marginal{\uncertain{Précisément} \ill{} \ill{} $E$ \ill{} une structure \ill{}
\ill{} \ill{} \ill{} \ill{} \ill{} \ill{} $E$ \ill{}}
\note{note écrite en diagonale dans la marge gauche, en bas du feuillet ;
presque entièrement illisible}

\page{168}
Je dis que cette loi de composition fait de $E$ un \emph{schéma en
groupes}. Il faut vérifier

(i) L'associativité, qui s'exprime par l'égalité de \emph{deux morphismes}
$E \times E \times E \to E$ au dessus de $G \times G \times G \to G$,
compatibles avec $\Gamma \times \Gamma \times \Gamma \to \Gamma$ ; cela
résulte du fait qu'il y a \emph{un seul} tel morphisme, à translation près
par un élément de $\underline{\mathrm{Hom}}_S(G \times G \times G, \Gamma)
\simeq \Gamma(S)$, donc un seul pour lequel l'unité va en l'unité

(ii) $e$ est l'unité, \uncertain{même} argument, \struck{\ill{}} \ill{}
\uncertain{appliqué} deux \ill{} \struck{\ill{}} sur l'identité

(iii) Existence d'une symétrie dans $E$.

Voir l'exposé du livre de Serre \struck{\ill{}} p.\ 183.

----

\marginal{N.B. Ce résultat est valable, de façon indépendante des
\uncertain{prop.} plus haut, si $G$ \ill{} \ill{} à fibres connexes
\ill{} ...}
\note{écrit en oblique dans la marge gauche}

Application au cas où $\Gamma$ \ill{} \ill{} \emph{étale},
\marginal{et $G$ à fibres géom. connexes}
\emph{Alors}
\[
\mathrm{Ext}^1_{S\text{-gr}}(G, \Gamma) \simeq \mathrm{Ext}^1_{S_0\text{-gr}}(G_0, \Gamma_0)
\]
où $S_0 = S_{\mathrm{réd}}$, $G_0 = G \times_S S_0$,
$\Gamma_0 = \Gamma \times_S S_0$. Ceci vaut plus gén. si $S_0 \to S$ est un
homéom. universel \uncertain{quelconque}. \struck{Ceci permet l'\ill{}}
\note{« et $G$ à fibres géom. connexes » est dans une bulle au-dessus de la
ligne}

\page{169}
\struck{\ill{} \ill{} multiplication \ill{} \ill{}, p.\ \ill{}}

D'ailleurs \struck{on} \ill{} une extension de $G$ par $\Gamma$ \ill{} par
la \uncertain{méthode} \ill{} \ill{}, en appliquant à
$\mathrm{Ext}^1_{S\text{-gr}}(G,\Gamma)$ les méthodes de descente [ainsi,
pour $S'$ \ill{} \add{sur $S$}, on obtient une \uncertain{fonction} en $S'$
compatible avec la descente fid. plate ....]

Prenons, p.\ ex.\ $G = \mathbb{G}_a$, $\Gamma = \mathbb{Z}/p\mathbb{Z}$,
\ill{} (\ill{} \ill{} \ill{} l'hypothèse noethérienne)
\[
\mathrm{Ext}^1_{S\text{-gr}}(\mathbb{G}_a, \mathbb{Z}/p\mathbb{Z}) \simeq \Gamma(S, \mathcal{O}^{p^{-\infty}})
\]
où
\[
\mathcal{O}^{p^{-\infty}} = \varinjlim_n (\mathcal{O}, n), \ldots
\]
l'hom.\ $(\mathcal{O}, n) \to (\mathcal{O}, m)$ défini par
$x \mapsto x^{p^{m-n}}$. On peut donc dire que le foncteur
$\underline{\underline{\mathrm{Ext}}}^1_{S\text{-gr}}(\mathbb{G}_a, \mathbb{Z}/p\mathbb{Z})$
est représentable \struck{par} \add{dans la catégorie des} \struck{les}
pro-schémas \uncertain{parfaits} sur $S$....

\marginal{\uncertain{Pour la dim.}, cf exposé Groth. \ill{} \ill{} sur les
\uncertain{isogénies} (Exposé 7)}
\note{écrit en oblique dans la marge gauche, en face de la formule}

\section*{Dualité de Cartier des groupes affines et groupes formels}
\note{inscrit à l'encre sur un feuillet de garde orange (page 170), qui ne
porte rien d'autre ; le feuillet ne reçoit pas de numéro de page. Ce titre
est le sien}

\page{171}
\note{tableau à double entrée, $V$ en lignes, $F$ en colonnes (la case
d'angle porte « $V$ / $F$ » séparés par une diagonale) ; chaque case
contient des groupes entourés d'un ovale au crayon}

\begin{itemize}
\item $V$ isom : $F$ isom : $(\mathbb{Z}/\ell)$ ; $F$ épim :
$(\mathbb{G}_m, \mathbb{Z}/\ell)$ ; $F$ nilpotent : $(\mu_p)$.
\item $V$ \struck{\ill{}} \uncertain{mixte} ? : $F$ isom : discret sép.\
\ill{} \ill{}, $(\mathbb{Z}/\ell, \mathbb{Z}/\ell$\struck{$\mathbb{Z}$}$)$ ;
$F$ épim : \struck{mult.\ \ill{} \ill{} \ill{} \ill{} \ill{} discret \ill{}}
$(\mathbb{G}_m, \mathbb{Z}/\ell\mathbb{Z}, \mathbb{Z})$ ; $F$ nilpotent :
$(\mu_p, \mathbb{V})$.
\item $V$ nilpotent : $F$ isom : $(\mathbb{Z}/p\mathbb{Z})$ ; $F$ épim :
$(\mathbb{G}_a, \mathbb{Z}/p\mathbb{Z})$ ; $F$ nilpotent :
\struck{\ill{} \ill{} radiciel} $(\alpha_p)$.
\end{itemize}
\note{les en-têtes des lignes et colonnes sont « isom », « épim »,
« nilpotent » ; l'en-tête de la deuxième ligne est un mot biffé suivi d'un
mot court \ill{}, que l'ordre des colonnes invite à lire « épim »}

\ill{} propriétés de $F$, $V$, caractérisation \ill{} l'existence de
\ill{}, de \uncertain{composantes} \ill{} des propriétés données ....

\page{172}
\begin{itemize}
\item $\mathbb{G}_m$ : $F$ épim, $V$ isom
\item $\mathbb{Z}/\ell\mathbb{Z}$ : $F$, $V$ isom
\item \struck{$\mu$}$\mu_p$ : $F = 0$, $V$ isom
\item $(\mathbb{G}_a,\ \alpha_p,\ \mathbb{V})$ : $F$ épim, $V = 0$ ; $F = 0$,
$V = 0$ ; $F = 0$, $V$ \struck{\ill{}} \uncertain{épim}
\item $\mathbb{Z}/p\mathbb{Z}$ : $F$ isom, $V = 0$ \quad $\ast$
\item $\mathbb{Z}/\ell\mathbb{Z}$ : $F$, $V$ isom
\item $\mathbb{Z}$ : $F$ isom, $V$ \uncertain{isom}
\item $Ab$ : $F$ épim, $V$ n'importe quoi
\end{itemize}
\note{une ligne d'en-têtes, les conditions sous chacun ; $Ab$ est séparé
des autres par un trait vertical ; les trois colonnes $\mathbb{G}_a$,
$\alpha_p$, $\mathbb{V}$ sont réunies par un crochet}

\begin{itemize}
\item $F$ isom : groupe discret séparable
\item $F^n = 0$ : groupe infinitésimal
\item $F$ épim : gr.\ \add{\ill{}} \struck{\ill{}} séparable (peut avoir
une partie $\mathbb{Z}$, mais pas de facteur de composition $\mathbb{V}$),
$\overset{?}{\Longrightarrow}$ i.e.\ pas de partie infinitésimale
\end{itemize}

\begin{itemize}
\item $V$ isom : groupe multiplicatif (\ill{})
\item $V^n = 0$ : groupe \struck{\ill{}} unipotent
\item $V$ \struck{épim} : \struck{groupe infinit.} pas de partie
\add{unipotente} \ill{} $\overset{?}{\Longrightarrow}$ en facteur de
composition.
\end{itemize}
\note{deux colonnes, $F$ à gauche, $V$ à droite}

[Condition de finitude \ill{} et symétrie \ill{} un catalogue \ill{} \ill{}
de groupe \ill{} \ill{} \ill{} \ill{} \ill{}]

[groupe infinitésimal $\neq$ groupe radiciel ;
$\mathbb{V}$ n'est pas unipotent ? bien qu'il soit une extension
\uncertain{d'unipotents}.]

$F$ \struck{\ill{}} $\Rightarrow$ $F$ isom ; $V$ épim $\Rightarrow$ $V$ isom (?)
\note{les deux implications sont réunies par une accolade suivie d'un point
d'interrogation}

Il y a 8 types fondamentaux de groupes affines (satisfaisant la condition
de finitude). Le groupe \ill{} \ill{} induction \uncertain{resp.} \ill{}
\add{unipotent} \ill{} types (les types de Dieudonné). Je pense que les
tableaux ci-dessus seront utiles \uncertain{incessamment}.

La \uncertain{première} \ill{} opposition \ill{} types de dualité, \ill{} :
dualité des var.\ abéliennes, \uncertain{d'une part}, \ill{} \ill{} des
dualités \ill{} \ill{}
\note{une flèche relie ce dernier paragraphe (à gauche) au paragraphe
précédent (à droite)}

\page{173}
\begin{itemize}
\item $\mathbb{G}_m$ : $F$ \uncertain{épim, non isom} ; $V$ isom
\item \struck{$\mu_p$} : $F = 0$ ; $V$ isom
\item $\mathbb{Z}/\ell\mathbb{Z}$ : $F$, $V$ isom
\end{itemize}
\note{le reste du feuillet est vierge}

\page{174}
\note{copie très pâle du feuillet suivant (page 175) : mêmes quatre
schémas, même N.B., mêmes listes accolées ; non retranscrite, voir page
175}

\page{175}
\note{quatre schémas disposés en deux colonnes séparées par un double
trait ; chacun a un titre en marge, des accolades nommant les parties, et
des mentions en oblique}

multiplicatif (\uncertain{$k$-simple connexe}) :

\begin{tikzcd}[column sep=small, row sep=small]
 & & \mathbb{Z}/\ell\mathbb{Z} \\
\mathbb{G}_m \arrow[r] & G \arrow[ur] \arrow[dr] & \\
 & & \mu_p
\end{tikzcd}

\note{accolade supérieure ($\mathbb{G}_m$ à $\mathbb{Z}/\ell\mathbb{Z}$) :
« \struck{\ill{}} séparable » ; inférieure ($\mathbb{G}_m$ à $\mu_p$) :
« connexe » ; à droite, « fini \struck{\ill{}}séparable » pour
$\mathbb{Z}/\ell\mathbb{Z}$, « fini radiciel » pour $\mu_p$, et une
accolade « fini » pour les deux}

discret (ordinaire) :

\begin{tikzcd}[column sep=small, row sep=small]
 & & \mathbb{Z}/\ell\mathbb{Z} \arrow[dl] \\
\mathbb{Z} & G \arrow[l] & \\
 & & \mathbb{Z}/p\mathbb{Z} \arrow[ul]
\end{tikzcd}

\note{accolade supérieure : « sans $p$-torsion » ; inférieure :
« \struck{\ill{}} sans $\ell$-torsion » ; devant $\mathbb{Z}$ :
« \struck{\ill{}} sans torsion » ; « fini multiplicatif » pour
$\mathbb{Z}/\ell\mathbb{Z}$, « fini unipotent » pour
$\mathbb{Z}/p\mathbb{Z}$, accolade « fini »}

unipotent (\uncertain{$k$-simple connexe}) :

\begin{tikzcd}[column sep=small, row sep=small]
 & & \mathbb{Z}/p\mathbb{Z} \\
\mathbb{G}_a \arrow[r] & G \arrow[ur] \arrow[dr] & \\
 & & \alpha_p
\end{tikzcd}

\note{accolade supérieure : « \struck{\ill{}}séparable » ; inférieure :
« connexe » ; « fini séparable » pour $\mathbb{Z}/p\mathbb{Z}$, « fini
radiciel » pour $\alpha_p$, accolade « fini »}

connexe (= \uncertain{infinitésimal} ?) :

\begin{tikzcd}[column sep=small, row sep=small]
 & & \mu_p \arrow[dl] \\
\mathbb{V} & G \arrow[l] & \\
 & & \alpha_p \arrow[ul]
\end{tikzcd}

\note{accolade supérieure : « sans torsion unipotente » ; inférieure :
« sans torsion multiplicative » ; devant $\mathbb{V}$ :
« \struck{\ill{}} sans torsion » ; « fini multiplicatif » pour $\mu_p$,
« fini unipotent » pour $\alpha_p$, accolade « fini »}

\emph{N.B.} Dans le \uncertain{cadre} des catégories, des
\struck{\ill{}} groupes \add{\uncertain{formels}} \ill{} des groupes
\ill{}, \ill{} les \ill{} \ill{}, le multiplicatif et l'additif
\uncertain{se mélangent} (dans le fini)

torsion $=$ \struck{\ill{} \ill{} \ill{} \ill{} \ill{}} $=$ \ill{}
sous-groupe fini
\note{à droite du N.B., un bloc encadré et barré de hachures, dont on ne
lit que ces mots}

\begin{itemize}
\item multiplicatif --- discret ordinaire
\item unipotent --- purement infinitésimal
\end{itemize}

\begin{itemize}
\item connexe séparable --- \struck{séparable} sans torsion
\item fini séparable --- fini multiplicatif
\item fini radiciel --- fini unipotent
\end{itemize}

\begin{itemize}
\item \struck{fini séparable} connexe --- sans torsion multiplicative
\item séparable --- sans torsion additive
\item fini --- fini
\end{itemize}

\begin{itemize}
\item multiplicatif connexe séparable --- gr. discret sans torsion
\item unipotent connexe séparable --- groupe infinitésimal
\struck{séparable} sans torsion \marginal{(= \struck{\ill{}} \ill{})}
\item fini séparable mult. \struck{connexe séparable} $=$ fini sép.
multiplicatif
\item fini radiciel mult. $=$ fini sép. unipotent
\item fini radiciel unipotent $=$ fini radiciel unipotent
\end{itemize}
\note{chaque groupe de lignes est réuni à gauche par une accolade ; la
note marginale, entourée de boucles, est presque entièrement biffée}

\page{176}
\[
\begin{array}{c|c|c|c}
\text{multiplicatif} & \mathbb{G}_m & \mathbb{Z}/\ell\mathbb{Z} & \mu_p \\ \hline
\text{unipotent} & \mathbb{G}_a & \mathbb{Z}/p\mathbb{Z} & \alpha_p
\end{array}
\]
\note{deux accolades sous le tableau : « $k$-simples » sous les colonnes
$\mathbb{G}_m$ et $\mathbb{Z}/\ell\mathbb{Z}$, « finis » sous les colonnes
$\mathbb{Z}/\ell\mathbb{Z}$ et $\mu_p$ ; le reste du feuillet est vierge}

\page{177}
\note{en haut, deux figures d'ovales qui se chevauchent, séparées par un
trait vertical ; les groupes sont inscrits dans les régions et des traits
obliques (certains repassés, d'autres biffés) les séparent. On donne les
étiquettes des ovales et les groupes lisibles}

À gauche : ovales étiquetés « $k$-connexe », « multiplicatif \ill{} »,
« fini séparable », « unipotent » (en bas à droite, souligné), « multiplicatif
séparable », « multiplicatif » (en bas, souligné), « \uncertain{unipotent}
$k$-connexe » ; groupes : $(\mathbb{G}_m)$, $(\mu_p)$,
$(\mathbb{Z}/\ell\mathbb{Z})$, $(\mathbb{G}_a, \alpha_p)$,
$(\mathbb{Z}/p\mathbb{Z})$, $(\mathbb{Z}/\ell\mathbb{Z})$,
\struck{$(\mathbb{Z}/p\mathbb{Z})$}, $(\mu_p)$,
$(\mathbb{G}_a, \alpha_p)$, et un groupe noirci \ill{}.
\note{« \struck{\ill{}} » biffé dans l'ovale supérieur}

À droite : ovales étiquetés « infinitésimal », « discret sép. »,
« sans torsion additive » (avec « \struck{infinitésimal} »), « discret
\struck{(séparable)} sans torsion multiplicative » ; groupes :
$(\alpha_p, \mathbb{V})$, $(\mu_p)$, $\mathbb{Z}/\ell\mathbb{Z}$,
\struck{$\mathbb{Z}/p\mathbb{Z}$}, $(\mathbb{Z}/\ell\mathbb{Z})$,
$(\mathbb{Z}/p\mathbb{Z})$, $(\mathbb{Z})$, $(\mu_p, \mathbb{V})$,
$(\mu_p)$.

Supposons la partie \struck{réduite} séparable \uncertain{unipotente}
connexe \uncertain{déf./$k$}. Dans le diagramme \ill{} \ill{} complété :

\note{un diagramme en losange : $(\mathbb{G}_m)$, $(\mu_p)$,
$(\mathbb{Z}/\ell\mathbb{Z})$, $(\mathbb{G}_a)$, $(\alpha_p)$,
$(\mathbb{Z}/\ell\mathbb{Z})$, $(\mathbb{Z}/p\mathbb{Z})$,
$(\mathbb{Z}/\ell\mathbb{Z})$, $(\mathbb{Z}/p\mathbb{Z})$, $(\alpha_p)$,
$(\mathbb{Z}/\ell\mathbb{Z})$, $(\mu_p)$, séparés par des traits obliques,
un ovale « fini » en haut}

Si de plus la partie séparable de $G_0$ est définie sur $k$, \ill{}
\struck{\ill{}} \ill{} \ill{} \ill{} de $G_0$ \ill{} \ill{}) dans \ill{}
un diagramme \struck{\ill{}} plus complet :

\[
\mu_p,\ (\mathbb{G}_a, \alpha_p),\ \mathbb{Z}/p\mathbb{Z},\ \mu_p,\ (\alpha_p, \mathbb{V}),\ \mathbb{Z}/p\mathbb{Z}
\]
\note{accolades : les trois premiers, les trois derniers, puis l'ensemble,
avec la mention « ext. \ill{} \ill{} »}
\[
\mathbb{G}_m \quad \bigl[\,(\mu_p),\ (\mathbb{G}_a, \alpha_p, \mathbb{V}),\ (\mathbb{Z}/p\mathbb{Z})\,\bigr] \quad (\mathbb{Z})
\]
\note{sous le crochet, une accolade « Ext. \ill{} \ill{} » réunit
$(\mathbb{G}_a, \alpha_p, \mathbb{V})$ ; en dessous, séparés par des
traits, $(\mathbb{Z}/\ell\mathbb{Z})$, $\mathbb{Z}/\ell\mathbb{Z}$,
$(\mathbb{G}_m)$, et un crochet vide ; un $\mu_p$ biffé à gauche}

\emph{Multiplicatif} \struck{\ill{}} \emph{connexe} (\uncertain{produit})
$\bigl($non-$p$-primaire $\times$ gr.\ \ill{}$\bigr)$ \ \emph{\ill{}
\ill{}} \emph{sans torsion}

\page{178}
\note{copie très pâle du feuillet suivant (page 179) : mêmes formules et
même diagramme ; non retranscrite, voir page 179}

\page{179}
\[
\mu_p + (\mathbb{G}_a, \alpha_p) \qquad (\alpha_p, \mathbb{V}) + \mathbb{Z}/p\mathbb{Z}
\]
\[
\mu_p \quad \alpha_p \quad \mathbb{Z}/p
\]
\[
\mathbb{G}_a \quad \alpha_p \quad \mathbb{V}
\]
\note{sous cette ligne, une flèche vers la gauche}
\[
\mu_p \times \mathrm{Unip.}, \quad \mathbb{V}
\]

Multiplicative \quad unipotents \qquad\qquad discret ordinaire

\note{un diagramme de traits : « multiplicative connexe séparable »,
avec au-dessus « $G_0$ (\uncertain{réd.}) » et au-dessous « $G_0$ réd »,
« $G$ réd » ; puis, disposés en losange, « multiplicatif fini radiciel »,
« multiplicatif fini sép. », « mult. fini sép. », « multiplicatif radiciel
fini » ; un trait biffé à droite}

$G$ réd est un sous-groupe \ill{} si $G$ est multiplicatif, mais pas s'il
est unipotent.

$G_0$ est défini sur $k$ \struck{si $G$ est multiplicatif}
\struck{mais pas} [faut-il \uncertain{prendre} \uncertain{variétés}
\add{\ill{}}, \ill{} \ill{} \ill{} ???]

unipotent connexe (\ill{} \ill{} sép.) \qquad \emph{fini unipotent
séparable}

\struck{\ill{}} \quad
$(\mathbb{G}_m)$ \ $(\mu_p)(\mathbb{Z}/\ell\mathbb{Z})$ \ Unip.\ conn.
\note{à droite, un diagramme de traits : « Unip.\ conn. »,
$\mathbb{Z}/\ell\mathbb{Z}$, $\mathbb{Z}/p\mathbb{Z}$, $\mu_p$,
$\mathbb{Z}/\ell\mathbb{Z}$, « Unip.\ \ill{} », $\mathbb{G}_m$,
$\mathbb{Z}/\ell\mathbb{Z}$, $\mu_p$ ; en dessous, $(\mu_p)$ d'où partent
deux traits vers « mult.\ fini » et « connexe »}

\page{180}
Modules discrets \struck{\ill{}} de type fini sur
\[
A = \varprojlim_n \mathbb{W}_n[F, \mathbb{V}]/(V^n, FV - p)
\]
\note{un petit indice, peut-être $\sigma$, est accroché à $\mathbb{W}_n$ ;
le $\mathbb{V}$ entre crochets est repassé ; entre « $=$ » et
$\varprojlim$, un mot biffé \struck{\ill{}}}

(Catégorie $C_1$) \quad on considère $C_1^\circ = C_2$.

Et la sous-catégorie \add{pleine} $C_1'$ de $C_1^{*}$ formée des modules
localement annulés par une puissance de $F$, \add{et la sous-catégorie
$C_2'$} [\ill{} \ill{} $C_1^\circ$ \ill{} \ill{} \struck{\ill{}} des types
\uncertain{finis} \emph{additifs} unipotents].

La dualité de Cartier des groupes finis, définie par \uncertain{un} \ill{} de
type \uncertain{cocyclique}, définit une \emph{isomorphisme}
\[
C_1' \simeq C_2' .
\]
D'ailleurs, \ill{} objet $G$ de $C_2$ \ill{} \ill{} \ill{} \ill{} objet
\ill{} tel que $G/G'$ \struck{\ill{}} soit le dual d'un objet quelconque de
$C_2'$, et \ill{} $=$ \ill{} objet à $C$, \ill{} dual \ill{} un objet
\struck{\ill{}} \ill{} de $C_1'$. Un objet de la \ill{} catégorie $C$ serait
un \struck{couple} \add{triple} formé de $(X, Y, u)$, où $X \in C_2$,
$Y \in C_1$, et $u$ un isom.\ de $\alpha(X)$ sur $\beta(Y)$

\begin{tikzcd}[row sep=small]
C_1 \arrow[dr, "\alpha"] & \\
 & C_0 \\
C_2 \arrow[ur, "\beta"'] &
\end{tikzcd}

\note{le diagramme donne $\alpha$ sur $C_1$ et $\beta$ sur $C_2$, alors que
la phrase écrit $\alpha(X)$ pour $X \in C_2$ et $\beta(Y)$ pour $Y \in
C_1$ ; transcrit tel quel. Dans la marge gauche, en face de « La dualité de
Cartier », « $C_1'$ $\simeq$ $C_2'$ » écrit verticalement, et les
étiquettes $C_1$ \struck{Modules de} et $C_2$}

\end{document}
